% sections/chapter-methodology.tex
\chapter{Methodology}
\label{ch:methodology}

\section{Field Sites}

The study draws on eighteen months of fieldwork (April 2024 -- September
2025) across three organisations, each anonymised per the agreement
described in the Preface. Table~\ref{tab:sites} summarises the sites and
the instrumentation available at each.

\begin{table}[htbp]
\centering
\small
\begin{tabular}{@{}p{2.1cm}p{2.6cm}p{2.6cm}p{4.4cm}@{}}
\toprule
\textbf{Site} & \textbf{Team studied} & \textbf{Headcount} & \textbf{Instrumentation} \\
\midrule
Nexa\index{Nexa} & Platform research & 34 & Chat logs, calendar, ticket tracker \\
Synapse\index{Synapse} & ML applied research & 21 & Chat logs, calendar, wiki edit history \\
Vertex\index{Vertex} & Infrastructure engineering & 46 & Chat logs, calendar, ticket tracker, incident log \\
\bottomrule
\end{tabular}
\caption{Field sites and available instrumentation.}
\label{tab:sites}
\end{table}

All three organisations run fully distributed or majority-remote research
teams spanning at least three time zones, and all three had, independently
of the study, already adopted some form of asynchronous-first
communication norm -- a precondition for the site-selection strategy, since
the phenomenon under study requires exactly this kind of environment to
appear at measurable scale.

\section{Data and Coding}

Three data streams were combined for each participant who consented to
instrumentation review: (1)~timestamped message and thread metadata from
the team's chat platform, stripped of message content and retained only as
channel-switch events; (2)~calendar data, coded for meeting type and
attendance; and (3)~a twice-weekly experience-sampling survey delivered
during working hours, asking participants to characterise their current
activity along a four-point scale from focused task work to active
reconciliation of competing status information.

Four hundred and twelve hours of meeting transcript were separately coded
by two research assistants for coordination-talk -- utterances whose
function was to establish or confirm ownership, status, or next steps
rather than to advance the substantive task at hand.\footnote{Inter-rater
reliability on the coordination-talk category reached
Cohen's~$\kappa = 0.81$ across a double-coded 15\% sample, which we treat as
adequate given the inherent ambiguity of utterances that do double duty.}
Coding used a scheme adapted from prior work on distributed cognition in
operational settings \citep{hollan2000distributed}, extended with two new
categories -- \emph{ownership drift} and \emph{artefact reconciliation} --
that emerged inductively from the first month of pilot coding at
Nexa\index{Nexa}.

\section{Constructing the Ambient Coordination Load Index}

The Ambient Coordination Load Index (ACLI)\index{Ambient Coordination Load
Index (ACLI)} combines three components for each participant-day:
\begin{enumerate}[label=(\roman*)]
  \item the count of distinct communication channels touched;
  \item the count of channel switches occurring within a 15-minute window
        of an unrelated channel touch, a proxy for reconciliation activity;
  \item self-reported reconciliation activity from the experience-sampling
        survey, weighted at 0.4 relative to the behavioural components.
\end{enumerate}
The index was validated against the coded transcript data by regressing
coder-identified reconciliation utterances per meeting against the
ACLI computed for that meeting's attendees in the surrounding hour
($R^2 = 0.58$, Chapter~\ref{ch:findings}). This is a moderate but, we
argue, meaningful correlation given that the transcript coding captures
only synchronous reconciliation, while the ACLI is designed to capture
asynchronous reconciliation as well.

\begin{editorialnote}
The full ACLI computation script and the de-identified participant-day
dataset are available to reviewers on request through the Press's
double-blind data-access portal; they are withheld from this manuscript
per the anonymised-review policy.
\end{editorialnote}

\section{Limitations}

Three limitations bear directly on how the findings in
Chapter~\ref{ch:findings} should be read. First, all three sites are
software-adjacent research organisations with above-average tooling
maturity; the ACLI's channel-count component may behave differently in
settings with fewer, more consolidated communication tools. Second, the
experience-sampling survey's four-point scale is a coarse instrument, and
the 0.4 weighting of self-report was chosen to balance reliability against
sensitivity rather than derived from an independent validation study.
Third, and most importantly, the eighteen-month window captures each
organisation at a single stage of its coordination-tooling maturity; we
cannot speak to how ACLI evolves as a team's tool stack consolidates or
fragments over a longer horizon.
